\scanpage{66}

et la condition que les $u_n$ définissent bien un hom de complexes est la ``condition de Stokes''
\[
K_n\partial_{n+1}=K_{n+1}d:
\]
\diagram{\begin{tikzcd}
A_{n+1}^n\arrow[r,"\partial_{n+1}"]\arrow[d,"d^n"']&
A_n^n\arrow[d,"K_n"]\\
A_{n+1}^{n+1}\arrow[r,"K_{n+1}"']&M
\end{tikzcd}}

Dans ces conditions, pour tt $X_*$\inserted{$\in\operatorname{Ob}\widehat{\Delta}$}, on a un hom. \unsure{can.} déduit de $u$
\diagram{\begin{tikzcd}
\Gamma(X_*,A_{X_*}^*)\arrow[r]&C^*(X_*,M_*)\arrow[d,equal]\\
&\Gamma(X_*,C_{\unsure{*}M}^*|X_*)
\end{tikzcd}}
qui représente l'hom dans la cat. dérivée
\[
\Gamma(X_*,A_{X_*}^*)\mathbin{-}\mathbb R\Gamma_{X_*}(M)
\]
associé à $u_*$. On aura alors \unsure{commutativité} dans
\diagram{\begin{tikzcd}[column sep=small]
\deleted{\mathbb R}\Gamma(X_*,A_{X_*}^*)\arrow[r]\arrow[d]&
C^*(X_*,M)=\Gamma(X_*,C_M^*|\unsure{X_*})\arrow[d,"\sim"]\\
\mathbb R\Gamma(X_*,A_{X_*}^*)\arrow[r]&
\mathbb R\Gamma(X_*,M)
\end{tikzcd}}

Si $A^*$ est une résolution de $M_*$, et $K_0:A_0^0\longrightarrow M$ induit l'identité sur $M\subset A_0^0$, alors la dernière flèche horizontale est aussi un iso, et

\margin{65}
